% !TEX root = ../main.tex

\chapter{绪论}

\section{研究背景与临床意义}
  前交叉韧带（Anterior Cruciate Ligament, ACL）撕裂是最常见的运动损伤之一。
  流行病学研究显示，此类损伤高度集中于15至30岁的年轻活跃人群\cite{sandersIncidenceAnteriorCruciate2016}，
  且常伴随半月板及软骨的继发性损伤\cite{spindlerClinicalPracticeAnterior2008}。
  相比之下，多发韧带膝关节损伤（Multi-Ligament Knee Injury, MLKI）虽然仅占膝关节韧带损伤的一小部分（约2.5\%）\cite{wilsonEpidemiologyMultiligamentKnee2014}，
  但其致伤原因多为车祸等高能创伤，往往导致严重的关节失稳，需要更为复杂的手术重建干预\cite{moatsheDiagnosisTreatmentMultiligament2017}。
  这两种损伤会在急性期影响患者运动功能，也会增加创伤后骨关节炎（Post-traumatic Osteoarthritis, PTOA）风险\cite{lohmanderLongtermConsequenceAnterior2007, snoekerRiskKneeOsteoarthritis2020}。
  \par 在疾病诊断与预后预测中，磁共振成像（Magnetic Resonance Imaging, MRI）凭借较好的软组织对比度，已成为评估膝关节状态的重要临床工具。
  基于三维 MRI 分割提取的定量指标，例如软骨三维厚度分布\cite{ecksteinMagneticResonanceImaging2006}和骨髓病变（Bone Marrow Lesions, BML）体积\cite{felsonAssociationBoneMarrow2001}，可用于观察骨关节炎（Osteoarthritis, OA）早期改变、评估治疗效果和随访患者结局。
  \par 然而，要实现这些高维指标的自动化提取，仍面临明显的临床挑战。
  首先是手动标注的时间成本较高：膝关节包含细小且复杂的软骨曲面和弥漫性 BML，
  放射科专家手动勾画单例三维高分辨率 MRI 往往需要耗费数小时，且存在显著的观察者间变异性。
  其次，当前基于深度学习的监督模型泛化性仍有限：不同医院、不同扫描仪及不同序列（如双回波稳态（Double Echo Steady-State, DESS）与临床脂肪抑制快速自旋回波（Fat-suppressed Fast Spin Echo, FSE）序列）之间存在较大的数据分布差异（Domain Shift）。
  由于真实的临床 FSE 大样本数据通常缺乏专家标注，传统的监督学习网络在此类目标域数据上可能出现明显的泛化性能下降。
  因此，有必要建立一种尽量减少目标域标注依赖、同时缓解跨序列域偏移的自动分割与量化流程。

\section{医学图像分割与无监督域适应技术现状}
  随着深度学习（Deep Learning）在计算机视觉领域的飞速发展，基于卷积神经网络（如U-Net、V-Net及其变体）的医学图像分割技术已在学术界确立了性能基准。
  针对临床数据缺少标签、不同序列间分布差异明显的问题，无监督域适应（Unsupervised Domain Adaptation, UDA）成为医学图像分割中的一个可行方向。
  UDA利用带有标签的源域数据（如公开的骨关节炎倡议（Osteoarthritis Initiative, OAI）数据集）训练模型，并通过特征级对齐（如域对抗神经网络（Domain-Adversarial Neural Network, DANN））或图像级翻译（如CycleGAN），使模型在无标签目标域上保持一定分割能力。

\section{本文的主要工作}
  本文面向临床大样本无标注膝关节 FSE 序列，建立骨与软骨无监督域适应分割和三维定量分析流程。模型训练采用无监督域适应策略，目标域人工标注仅用于小规模独立测试集验证，不参与训练。主要工作如下：
  \subsection{提出基于2D U-Net改进的域适应分割模型}
    针对跨序列分布差异，提出了一种改进的 DANN 域对抗模型。
    通过引入实例归一化（Instance Normalization, IN）尝试减弱底层对比度风格差异，并结合空洞空间金字塔池化（Atrous Spatial Pyramid Pooling, ASPP）设计软骨多尺度上下文建模模块。
    同时，设计基于损失比率的动态冻结策略，以缓解判别器收敛过快带来的训练失衡，并在人工标注测试集上验证该模型用于 FSE 序列骨与软骨自动提取的可行性。
  \subsection{通过软骨厚度和 BML 候选高信号体积等临床指标对比 ACL 和 MLKI 的损伤模式差异}
    开发了基于欧氏距离变换的软骨厚度映射、基于kd-tree的胫股关节影像近接触面积计算工具，以及面向FSE高信号候选区域的BML候选高信号体积量化流程。
    并对质量控制后纳入的430例孤立ACL撕裂与149例MLKI患者的客观提取数据进行了宏观统计检验。
    统计结果显示，两类损伤在胫股关节影像近接触面积与自动提取的BML候选高信号体积分布上存在差异，
    可为后续研究两类损伤的影像学差异和病理演进提供定量参考。

\section{论文结构安排}
  本文共分为六章。
  第一章介绍研究背景与意义；
  第二章综述相关领域的文献与技术现状；
  第三章详述跨域数据集的构建与标准化预处理；
  第四章重点论述改进型DANN域对抗网络的设计与单平面实验结果；
  第五章阐述自动量化流程的开发，并进行ACL与MLKI的临床统计对比；
  第六章为全文总结与展望。
  \par 本文对不同组织的验证条件并不相同：骨与软骨具备人工测试集，可进行Dice和表面距离等分割准确性评价；BML在本研究中缺乏大规模人工病灶金标准，因此文中的BML自动检测与分割主要指FSE高信号候选区域的自动检测、粗分割和体积量化，其精细病灶分割准确性仍需后续研究验证。
